\begin{abstract}
\noindent En una colocación legal de $n-1$ reinas sobre un tablero de $n\times n$, la única
casilla que puede recibir la última reina es la intersección de la fila y la columna
vacías. Estudiamos la probabilidad de que esa casilla esté libre de ataques cuando
la colocación se elige uniformemente. Cuando faltan $h\ge1$ reinas, las $h$ filas
y las $h$ columnas vacías se cortan en $h^2$ casillas, los cruces vacíos; sea $f_h(n)$ la fracción
media de ellos que admite una reina. Probamos que, para cada $h$
fijo, $f_h(n)\to F(0)=1/b_0>0$, donde $b_0$ es la constante de normalización del
perfil variacional en el tablero lleno. La prueba se apoya en la cota inferior de Simkin para el
número de soluciones y en la regularidad del perfil variacional hasta el tablero lleno, demostrada
en un trabajo complementario. Para $f_1$ y $f_2$ damos además una prueba directa de una cota
inferior asintótica explícita, aunque muy conservadora.

\smallskip\noindent Si $N_h(n)$ cuenta las colocaciones legales de $n-h$ reinas y $Q(n)=N_0(n)$,
se deduce que $N_h(n)/Q(n)\sim n^h/((h!)^2F(0)^h)$. Además, para cada $K$ fijo,
$\sum_{h\le K}\log[f_h(n)/F(h/n)]\to0$, donde $F(x)=(1-x)/b_x$ y $b_x$ es la constante de normalización
del perfil de masa $1-x$. La prueba
compara la ley uniforme con un proceso de referencia mediante entropía relativa
$o(n)$ y cuenta, en ambos sentidos, las reparaciones que conectan dos capas
consecutivas. Esta cuenta da también la ley media de la posición de los cruces vacíos y la
concentración del número de reparaciones por cruce.
\ifverificacionLean
Los enunciados principales también se han formalizado en Lean~4; al final del artículo se
precisa el alcance de esa verificación.
\fi
La estimación numérica $F(0)\approx0{,}0809$ carece de una cota de error certificada.

\medskip\noindent
\textbf{Palabras clave:} problema de las $n$ reinas; enumeración asintótica; completabilidad;
entropía relativa; absorbedores; método de ecuaciones diferenciales; doble cuenta.

\noindent
\textbf{MSC 2020:} 05A16; 05C70; 60C05; 94A17.
\end{abstract}
